\documentclass[12pt,a4paper]{ctexart}

% === 字体设置 ===
\setCJKmainfont{SimSun}
\setmainfont{Times New Roman}

% === 宏包 ===
\usepackage{amsmath}
\usepackage{amssymb}
\usepackage[version=4]{mhchem}
\usepackage{graphicx}
\usepackage{paracol}
\usepackage{xcolor}
\usepackage{fancyhdr}
\usepackage{geometry}

% === 页面设置 ===
\geometry{
  a4paper,
  left=1.5cm,
  right=1.5cm,
  top=2cm,
  bottom=2cm,
  columnsep=1cm
}

% === 内部辅助：计算 paracol 列宽 ===
\newlength{\colwidthinner}
\newcommand{\calccolwidth}{%
  \setlength{\colwidthinner}{\dimexpr\linewidth-2\fboxsep-2\fboxrule\relax}%
}

% 右栏：修改建议框
\newcommand{\correctionbox}[1]{%
  \calccolwidth
  \colorbox{blue!8}{\parbox{\colwidthinner}{#1}}%
}

% === 页眉页脚 ===
\pagestyle{fancy}
\fancyhf{}
\fancyhead[L]{校对报告}
\fancyhead[R]{\leftmark}
\fancyfoot[C]{\thepage}

\begin{document}

\title{第 1 讲校对测试1}
\author{}
\date{}
\maketitle

\section*{知识}
\begin{paracol}{2}

## 模块一 功和功率\\
\\
### 必备知识\includegraphics[width=\linewidth]{./images/image1.png}\{width="0.19444444444444445in" height="0.19444444444444445in"\}\\
\\
#### 一、判断力是否做功及做正、负功的方法\\
\\
判断依据                            适用情况\\
根据力和位移方向的夹角判断             常用于恒力做功的判断\\
\\
根据力和瞬时速度方向的夹角判断           常用于质点做曲线运动\textsuperscript{\textcolor{red}{\textcircled{3}}}\\
\\
根据功能关系或能量守恒定律判断           常用于变力做功的判断\\
\\
#### 二、计算功的方法\\
\\
#### **恒力做的功**\\
\\
直接用*$W=Fl\sin \alpha $*计算。\textsuperscript{\textcolor{red}{\textcircled{1}}}\\
\\
#### **合外力做的功**\\
\\
方法一：先求合外力*${F}_{合}$*，再用*${W}_{合}={F}_{合}l\cos \alpha $*求功。\textsuperscript{\textcolor{red}{\textcircled{4}}}\\
\\
方法二：先求各个力做的功*${W}_{1}$*、*${W}_{2}$*...，再应用：*${W}_{合}={W}_{1}+{W}_{2}+\...$*求合外力做的功。\textsuperscript{\textcolor{red}{\textcircled{5}}}\\
\\
#### **变力做的功**\\
\\
（1）应用动量定理求解；\textsuperscript{\textcolor{red}{\textcircled{2}}}\\
\\
（2）用$W=Pt$求解，其中变力的功率$P$不变；\\
\\
（3）常用方法还有转换法、微元法、图像法、平均力法等，求解时根据条件灵活选择。\\

\switchcolumn

\correctionbox{\textcircled{1} 恒力做功的定义为“力乘以力在位移方向上的投影与位移大小的乘积”，其中\$\textbackslash alpha\$是力与位移方向的夹角，应使用余弦分量，原文用正弦属于核心概念错误；同时LaTeX公式已默认对物理量使用斜体格式，无需额外添加斜体标记`*`，避免格式冗余。}
\medskip

\correctionbox{\textcircled{2} 动量定理描述的是合外力的冲量（力的时间积累）与动量变化的关系，无法直接用于求解做功（力的空间积累）；变力做功对应动能变化，应使用动能定理（合外力对物体做的总功等于物体动能的变化量），此处属于物理概念混淆。}
\medskip

\correctionbox{\textcircled{3} 表述不完整，与前后两条适用情况的表述结构不一致，明确该方法的适用场景为曲线运动下的做功判断，避免语义歧义。}
\medskip

\correctionbox{\textcircled{4} LaTeX公式中的物理量已默认使用斜体格式，无需额外添加斜体标记`*`，删除冗余格式标记，保证排版规范。}
\medskip

\correctionbox{\textcircled{5} 1. 公式外的斜体标记`*`冗余，予以删除；2. 中文语境下的省略号应使用全角“……”，公式内的省略号应使用规范的LaTeX命令`\textbackslash cdots`，避免使用不规范的半角“...”；3. 公式前无需额外添加冒号，修正标点使用错误。}
\medskip

\correctionbox{\textcircled{6} 当前配图显示为加载失败的占位符，属于图片缺失；且图片标识为“test”，无法确认内容与模块主题匹配，同时标题与图片同行排版不符合Markdown规范，应分行放置并替换为有效配图。

---}
\medskip

\switchcolumn*

\end{paracol}


ewpage

\section*{第1题}
\begin{paracol}{2}

**一本班1**（多选）\\
如图所示\textsuperscript{\textcolor{red}{\textcircled{1}}}，利用斜面从货车上卸货，每包货物的质量$m=20\mathrm{kg}$，斜面倾角$\alpha ={37}^{\circ }$，斜面的长度$l=0.5\mathrm{m}$，货物与斜面间的动摩擦因数$\mu =0.2$，货物从斜面顶端滑到底端的过程中，有（ ）（取$g=10\mathrm{m}/{\mathrm{s}}^{2}$）\textsuperscript{\textcolor{red}{\textcircled{2}}}\\
A.                                  重力对货物做的功为$60\mathrm{J}$\\
B.                                  斜面支持力对货物做的功为$0$\\
C.                                  斜面摩擦力对货物做的功为$16\mathrm{J}$\\
D.                                  合外力对货物做的功为$76\mathrm{J}$\\
答案:\\
AB\\
解答：                             重力做功：${W}_{\mathrm G}=mgL\sin {37}^{\circ }=60\mathrm{J}$\textsuperscript{\textcolor{red}{\textcircled{3}}}，选项$\mathrm{A}$正确；\\
支持力方向与位移方向垂直，不做功，则有：${W}_{\mathrm N}=0$，选项$\mathrm{B}$正确；\\
滑动摩擦力为：$f=\mu {F}_{\mathrm{N}}=0.2\times 20\times 10\times 0.8\mathrm{N}=32\mathrm{N}$，则摩擦力做功为：${W}_{\mathrm{f}}=-fL=-32\times 0.5\mathrm{J}=-16\mathrm{J}$\textsuperscript{\textcolor{red}{\textcircled{4}}}，选项$\mathrm{C}$错误；\\
合外力做功为：${W}_{合}={W}_{\mathrm{N}}+{W}_{\mathrm{f}}+{W}_{\mathrm{G}}=60\mathrm{J}-16\mathrm{J}=44\mathrm{J}$，选项$\mathrm{D}$错误。\\

\switchcolumn

\correctionbox{\textcircled{1} 题干提及“如图所示”但未提供对应配图，存在图片缺失，需补充配图以完善题干表述。}
\medskip

\correctionbox{\textcircled{2} 调整重力加速度说明的位置至问题前，补充多选题常规引导语，使题干表述通顺、逻辑清晰。}
\medskip

\correctionbox{\textcircled{3} 与题干中斜面长度的符号\$l\$保持一致，规范物理量符号使用。}
\medskip

\correctionbox{\textcircled{4} 与题干中斜面长度的符号\$l\$保持一致，规范物理量符号使用。

---}
\medskip

\switchcolumn*

\end{paracol}


ewpage

\section*{第2题}
\begin{paracol}{2}

**例1**\\
如图所示，重力为$G$的物体静止在倾角为$\alpha $的粗糙斜面体上，现使斜面体水平向左做匀速直线运动\textsuperscript{\textcolor{red}{\textcircled{1}}}，通过的位移为$x$，物体相对斜面体一直保持静止，则在这个过程中（ ）\includegraphics[width=\linewidth]{./images/image2.png}\{width="2.78125in" height="1.1666666666666667in"\}\\
A.                                  弹力对物体做功为$Gx\cos \alpha $\\
B.                                  静摩擦力对物体做功为$Gx\sin \alpha $\\
C.                                  重力对物体做功为零\\
D.                                  合力对物体做功为$Gx$\\
答案:\\
D\\
解答：                             $\mathrm{ABC}$．分析物体的受力情况：重力$mg$、弹力$N$和摩擦力$f$\textsuperscript{\textcolor{red}{\textcircled{5}}}，如图所示：\\
\includegraphics[width=\linewidth]{./images/image3.png}\{width="5.072916666666667in" height="4.0625in"\}\\
根据平衡条件，有：\\
$N=G\sin \theta $\\
$f=G\cos \theta $\\
重力与位移垂直，做功为零；\\
摩擦力$f$与位移的夹角为$\alpha $，所以摩擦力对物体$m$做功为\textsuperscript{\textcolor{red}{\textcircled{6}}}：\\
${W}_{f}=fx\cos \alpha =Gx\sin \alpha \cos \alpha $\\
斜面对物体的弹力做功为：\\
${W}_{N}=Nx\cos \left({90}^{\circ }+\alpha \right)=-Gx\sin \alpha \cos \alpha $；故$\mathrm{AB}$错误，$\mathrm{C}$正确；\\
$\mathrm{D}$．因物体做匀速运动，根据动能定理可知，合外力做功为零，故$\mathrm{D}$错误。\\
故选：$\mathrm{D}$。\textsuperscript{\textcolor{red}{\textcircled{4}}}\\

\switchcolumn

\correctionbox{\textcircled{1} 题干描述的运动方向与配图中速度箭头方向、位移方向（虚线斜面在实线右侧）矛盾，且与解答中摩擦力、弹力做功的夹角分析不符。}
\medskip

\correctionbox{\textcircled{2} 倾角符号应与题干的\$\textbackslash alpha\$统一；静止在斜面上的物体受力平衡，垂直斜面方向支持力等于重力的垂直分力\$G\textbackslash cos\textbackslash alpha\$，沿斜面方向静摩擦力等于重力的沿斜面向下分力\$G\textbackslash sin\textbackslash alpha\$，原表达式将两者写反。}
\medskip

\correctionbox{\textcircled{3} 重力方向与水平位移垂直，做功为零（选项C正确）；物体做匀速直线运动，动能变化为零，根据动能定理合力做功为零（选项D错误），正确答案为C。}
\medskip

\correctionbox{\textcircled{4} 解答分析已明确选项C正确、D错误，最终选择应与分析结论一致。}
\medskip

\correctionbox{\textcircled{5} 题干已明确物体重力为\$G\$，应统一使用该符号，避免引入未定义的质量符号\$m\$造成混淆。}
\medskip

\correctionbox{\textcircled{6} 题干未提及物体质量\$m\$，应删除多余的未定义符号，保持表述一致。}
\medskip

\switchcolumn*

\end{paracol}


ewpage

\section*{第3题}
\begin{paracol}{2}

**教师版**\\
我国古代劳动人民创造了璀璨的农耕文明。图（$\mathrm{a}$）为《天工关物》\textsuperscript{\textcolor{red}{\textcircled{1}}}中描绘的利用耕牛整理田地的场景，简化的物理模型如图（$\mathrm{b}$）所示，人站立的农具视为与水平地面平行的木板，两条绳子相互平行且垂直于木板边缘。已知绳子与水平地面夹角$\theta$为$25.5^{\circ}$，$\sin 25.5^{\circ}=0.43$，$\cos 25.5^{\circ}=0.90$。当每条绳子拉力$F$的大小为$250{\rm N}$时，人与木板沿直线匀速前进，在$15{\rm s}$内前进了$20{\rm m}$，求此过程中：\includegraphics[width=\linewidth]{./images/image4.png}\{width="2.8958333333333335in" height="1.5208333333333333in"\}\\
（1）地面对木板的阻力大小；\\
（2）两条绳子拉力所做的总功；\\
（3）两条绳子拉力的总功率。\\
解答：                              （1）对木板，水平方向受力平衡，则阻力$f=2F\cos \theta =2\times 250\times 0.9\mathrm{N}=500\mathrm{N}$\textsuperscript{\textcolor{red}{\textcircled{2}}}\\
（2）根据功的定义，拉力做功$W=2Fx\cos \theta =2\times 250\times 20\times 0.90\mathrm{J}=9000\mathrm{J}=9\times {10}^{4}\mathrm{J}$\textsuperscript{\textcolor{red}{\textcircled{4}}}\\
（3）根据功率的定义，拉力功率$P = \frac{W}{t} = \frac{9000}{15}\mathrm{W} = 600\mathrm{W}$\\
答：（1）地面对木板的阻力大小为$500\mathrm{N}$；\textsuperscript{\textcolor{red}{\textcircled{3}}}\\
（2）两条绳子拉力所做的总功为$9.0\times {10}^{4}\mathrm{J}$；\textsuperscript{\textcolor{red}{\textcircled{5}}}\\
（3）两条绳子拉力的总功率为$600\mathrm{W}$。
\switchcolumn

\correctionbox{\textcircled{1} 《天工开物》是明代宋应星所著古代科技著作的标准名称，此处“关”为错别字。}
\medskip

\correctionbox{\textcircled{2} 数值计算错误，2×250×0.9的正确结果为450，原结果不符合运算规则。}
\medskip

\correctionbox{\textcircled{3} 与（1）的正确计算结果对应，原答案数值错误。}
\medskip

\correctionbox{\textcircled{4} 科学计数法使用错误，9000对应的数量级为10³，原表述数量级错误。}
\medskip

\correctionbox{\textcircled{5} 与（2）的正确计算结果对应，原答案数量级错误。

---}
\medskip

\switchcolumn*

\end{paracol}

\end{document}
